%% BEAMER THEME FLIP 2012: Main tex file for compiling
%$ Compile this file. 
%%
%% Copyright 2012 by Flip Tanedo
%% This file may be distributed and/or modified
%% 	1. under the LaTeX Project Public License and/or
%% 	2. under the GNU Public License.
%% 
%% If you e-mail Flip (pt267@cornell.edu) to say that you
%% like this style file, then it would make him smile.

%% Please see notes.txt for comments on Beamer Theme Flip 2013
%% By default, this template is meant to be run with XeLaTeX (for fonts)
%% To run in PDFLaTeX, remove fontspec and any font commands

%% Discussion of Beamer vs XeLaTeX vs LuaLaTeX
%% http://tex.stackexchange.com/questions/29497/xelatex-preventing-beamer-from-using-different-backgrounds


\documentclass[12 pt]{beamer}
\usetheme[
	bullet=circle,		% Other option: square
	bigpagenumber,		% circled page number on lower right
	topline=true,			% colored bar at the top of the frame 
	shadow=false,			% Shading for beamer blocks
	watermark=BG_lower,	% png file for the watermark
	]{Flip}


\newcommand{\titleimage}{title}			% Custom title 
\newcommand{\tanedo}{tanedolight}		% Custom author name
\newcommand{\CMSSMDM}{CMSSMDMlight.png}	% light background plot


%%%%%%%%%%
% FONTS %
%%%%%%%%%%

%% Default font: lmodern, doesn't require fontspec % solves some default warnings
\usepackage[T1]{fontenc}
\usepackage{lmodern}			
%\usepackage{sfmath}		% Sans Serif Math, off by default

%% Protects fonts from Beamer screwing with them
%% http://tex.stackexchange.com/questions/10488/force-computer-modern-in-math-mode
\usefonttheme{professionalfonts}


%% XeLaTeX fonts: (comment out if you don't use XeLaTeX)

%% For advanced fonts: access local OS X fonts
%\usepackage[no-math]{fontspec}		
%% This template uses typical OS X and Adobe fonts
%\defaultfontfeatures{Mapping=tex-text}	% This seems to be important for mapping glyphs properly

%\setmainfont{Gill Sans}			% Beamer ignores "main font" in favor of sans font
%\setsansfont{Gill Sans}			% This is the font that beamer will use by default
%% \setmainfont{Gill Sans Light}		% Prettier, but harder to read

%\setbeamerfont{title}{family=\fontspec{Gill Sans}}


%\newcommand{\handwriting}{\fontspec{augie}} % From Emerald City, free font
\newcommand{\handwriting}{}	% If you prefer no special handwriting font or don't have augie

%% Gill Sans doesn't look very nice when boldfaced
%% This is a hack to use Helvetica instead
%% Usage: \textbf{\forbold some stuff}
%\newcommand{\forbold}{\fontspec{Helvetica}}
\newcommand{\forbold}{} % if you want no special boldface



%%%%%%%%%%%%%%%%%%%%%%%%
% Usual LaTeX Packages %
%%%%%%%%%%%%%%%%%%%%%%%%

\usepackage{amsmath}
\usepackage{amsfonts}
\usepackage{amssymb}
\usepackage{graphicx}
%\usepackage{mathrsfs} 			% For Weinberg-esque letters
\usepackage{cancel}				% For "SUSY-breaking" symbol
\usepackage{slashed}            % for slashed characters in math mode
%\usepackage{bbm}                % for \mathbbm{1} (unit matrix)
\usepackage{amsthm}				% For theorem environment
%\usepackage{multirow}			% For multi row cells in table
%\usepackage{arydshln} 			% For dashed lines in arrays and tables
%\usepackage{tikzfeynman}		% For Feynman diagrams
% \usepackage{subfig}           % for sub figures
% \usepackage{young}			% For Young Tableaux
% \usepackage{xspace}			% For spacing after commands
% \usepackage{wrapfig}			% for Text wrap around figures
% \usepackage{framed}


\graphicspath{{images/}{images-template/}}	% Put all images in this directory. Avoids clutter.


\usetikzlibrary{backgrounds}
\usetikzlibrary{mindmap,trees}	% For mind map
% http://www.texample.net/tikz/examples/computer-science-mindmap/


% SOME COMMANDS THAT I FIND HANDY
% \renewcommand{\tilde}{\widetilde} % dinky tildes look silly, dosn't work with fontspec
\newcommand{\comment}[1]{\textcolor{comment}{\footnotesize{#1}\normalsize}} % comment mild
\newcommand{\Comment}[1]{\textcolor{Comment}{\footnotesize{#1}\normalsize}} % comment bold
\newcommand{\COMMENT}[1]{\textcolor{COMMENT}{\footnotesize{#1}\normalsize}} % comment crazy bold
\newcommand{\Alert}[1]{\textcolor{Alert}{#1}} % louder alert
\newcommand{\ALERT}[1]{\textcolor{ALERT}{#1}} % loudest alert
%% "\alert" is already a beamer pre-defined



\author[Leonardo Maccari\quad {leonardo.maccari@unitn.it}]{Leonardo Maccari}
\title[Reti Comunitarie]{Reti Comunitarie}
\institute{DISI, University of Trento}
\date{\today}



\begin{document}

%%%%%%%%%%%%%%%%%%%%%%%%
% Additional  settings %
%%%%%%%%%%%%%%%%%%%%%%%%

%% To use external nodes; http://www.texample.net/tikz/examples/beamer-arrows/
\tikzstyle{every picture}+=[remember picture]


%%%%%%%%%%%%%%%%%%%%%%%%
% Title page           %
%%%%%%%%%%%%%%%%%%%%%%%%



% Macro per il titolo e altre variabili di report
\newcommand{\repnumber}{Secondo Report}


%misc
\newcommand{\un}{\uncover<+->}
\newenvironment{code}{\footnotesize\verbatim}{\endverbatim}
\newcommand{\dfixme}[2]{\fixme[margin]{#1}\fixme[footnote]{#2}}
% item e boxes

\newcommand{\bii}{\begin{itemize}[<+->]}
\newcommand{\bi}{\begin{itemize}}
\newcommand{\ei}{\end{itemize}}
\newcommand{\ii}{\item}

\newcommand{\bd}{\begin{description}}
\newcommand{\ed}{\end{description}}
\newcommand{\be}{\begin{enumerate}}
\newcommand{\ee}{\end{enumerate}}

\newcommand{\bab}[1]{\begin{alertblock}{#1}}
\newcommand{\eab}{\end{alertblock}}
\newcommand{\beb}[1]{\begin{exampleblock}{#1}}
\newcommand{\eeb}{\end{exampleblock}}
\newcommand{\bbox}{\begin{beamerboxesrounded}}
\newcommand{\ebox}{\end{beamerboxesrounded}}


% Standard e protocolli


\newcommand{\wi}{802.11}
\newcommand{\wm}{WiMax}
\newcommand{\fl}{ISM}

\newcommand{\std}{\ensuremath{\textit{Standard 802.16-2004}}\xspace} 
\newcommand{\pro}{\ensuremath{\textit{protocollo 802.16-2004}}\xspace}
\newcommand{\drae}{\ensuremath{\textit{Draft 802.16e}}\xspace} 
\newcommand{\proe}{\ensuremath{\textit{protocollo 802.16e}}\xspace}

\newcommand{\ui}{\ensuremath{\textit{IEEE 802.11i}}\xspace}
\newcommand{\us}{\ensuremath{\textit{IEEE 802.11}}\xspace}
\newcommand{\se}{\ensuremath{\textit{IEEE 802.16e}}\xspace}
\newcommand{\sd}{\ensuremath{\textit{IEEE 802.16d}}\xspace}
\newcommand{\Ss}{\ensuremath{\textit{IEEE 802.16}}\xspace}
\newcommand{\Sc}{\ensuremath{\textit{IEEE 802.16-2005}}\xspace}
\newcommand{\ux}{\ensuremath{\textit{IEEE 802.1X}}\xspace}

% sigle di nodi
\newcommand{\aut}{\ensuremath{\textit{Aut}}\xspace}
\newcommand{\aau}{\ensuremath{\textit{AA}_1}\xspace}
\newcommand{\aad}{\ensuremath{\textit{AA}_2}\xspace}
\newcommand{\aat}{\ensuremath{\textit{AA}_3}\xspace}
\newcommand{\aaq}{\ensuremath{\textit{AA}_4}\xspace}
\newcommand{\aaid}{\ensuremath{\textit{AA}_{id}}\xspace}
\newcommand{\stau}{\ensuremath{\textit{SA}_1}\xspace}
\newcommand{\stad}{\ensuremath{\textit{SA}_2}\xspace}
\newcommand{\sta}{\ensuremath{\textit{Sta}}\xspace}
\newcommand{\as}{\ensuremath{\textit{AS}}\xspace}

\newcommand{\sn}{\ensuremath{\textit{sponsor node}}\xspace}
\newcommand{\bs}{\ensuremath{\textit{base station}}\xspace}
\newcommand{\supp}{\ensuremath{\textit{supplicant}}\xspace}
\newcommand{\auth}{\ensuremath{\textit{authenticator}}\xspace}
\newcommand{\auths}{\ensuremath{\textit{authentication server}}\xspace}





% chiavi
\newcommand{\psa}{\textit{primary security association}}
\newcommand{\ssa}{\textit{static security association}}
\newcommand{\dsa}{\textit{dynamic security association}}
\newcommand{\fa}{\ensuremath{\textit{fast re-authentication}}\xspace}
\newcommand{\substa}{\ensuremath{\textit{SS}}\xspace}
\newcommand{\basesta}{\ensuremath{\textit{BS}}\xspace} 
\newcommand{\mobsta}{\ensuremath{\textit{MS}}\xspace}
\newcommand{\authserv}{\ensuremath{\textit{AS}}\xspace}  
\newcommand{\secass}{\ensuremath{\textit{SA}}\xspace}  
\newcommand{\authkey}{\ensuremath{\textit{AK}}\xspace} 
\newcommand{\tek}{\ensuremath{\textit{TEK}}\xspace}   

\newcommand{\pak}{\ensuremath{\textit{PAK}}\xspace}   
\newcommand{\ppak}{\ensuremath{\textit{Pre-PAK}}\xspace}   
\newcommand{\eik}{\ensuremath{\textit{EIK}}\xspace}   
\newcommand{\kek}{\ensuremath{\textit{KEK}}\xspace}   
\newcommand{\sa}{\ensuremath{\textit{SA}}\xspace}
\newcommand{\su}{\ensuremath{\textit{SU}}\xspace}
\newcommand{\pmk}{\ensuremath{\textit{PMK}}\xspace}
\newcommand{\psk}{\ensuremath{\textit{PSK}}\xspace}
\newcommand{\ptk}{\ensuremath{\textit{PTK}}\xspace}
\newcommand{\gtk}{\ensuremath{\textit{GTK}}\xspace}
\newcommand{\pmki}{\ensuremath{\textit{PMK}_i}\xspace}
\newcommand{\pmkid}{\ensuremath{\textit{PMKID}}\xspace}
\newcommand{\pmksa}{\ensuremath{\textit{PMKSA}}\xspace}
\newcommand{\ptksa}{\ensuremath{\textit{PTKSA}}\xspace}
\newcommand{\rnd}{\ensuremath{\textit{RANDOM}}\xspace}
\newcommand{\tkn}{\ensuremath{\textit{TOKEN}}\xspace}
\newcommand{\tknn}{\ensuremath{\textit{TOKEN}}}
\newcommand{\emsk}{\ensuremath{\textit{AMSK}}\xspace}
\newcommand{\amsk}{\ensuremath{\textit{AMSK}}\xspace}
\newcommand{\ak}{\ensuremath{\textit{AK}}\xspace}
\newcommand{\os}{\ensuremath{\textit{OSS}}\xspace}
\newcommand{\teko}{\ensuremath{\textit{TEK0}}\xspace}   
\newcommand{\teki}{\ensuremath{\textit{TEK1}}\xspace}   
\newcommand{\hmac}{\ensuremath{\textit{HMAC}}\xspace}   

\newcommand{\emskid}{\ensuremath{\textit{AMSKID}}\xspace}
\newcommand{\msk}{\ensuremath{\textit{MSK}}\xspace}
\newcommand{\radkey}{\ensuremath{\textit{RADIUSKEY}}\xspace}



{ %% This is a total kludge for a fancy title page background
\setbeamertemplate{sidebar right}{\llap{
\includegraphics[width=\paperwidth,height=\paperheight]{topo-trans-red}}}
\begin{frame}[c]%{\phantom{title page}} 
% The \phantom{title page} is a kludge to get the red bar on top
% \titlepage
\begin{center}
	% \includegraphics[width=7cm]{WarpedPenguinsReturn}

	\begin{tikzpicture}%[show background grid] %% Use grid for positioning, then turn off
		\node[inner sep=0pt,above right] (title) 
			{ \includegraphics[width=7cm]{title-dag} };
		% \node (title) at (1.5,1.5) {};
	\end{tikzpicture}
	\quad

	% \includegraphics[width=7cm]{\titleimage} 
	
	
	%\includegraphics[height=1.5cm]{\tanedo} \quad
	 % {\fontspec{Zapfino} Flip Tanedo} \quad
	% \includegraphics[height=1cm]{FlipSansSerif} \quad
	% \footnotesize\textcolor{gray}{In collaboration with} Csaba Cs\'aki\textcolor{gray}{,} Yuval Grossman\textcolor{gray}{, and} Yuhsin Tsai\normalsize\\
	% \textcolor{Comment}{ \;($\pi$ day)}\\
	% \Comment{4 February 2011}
 \footnotesize\textcolor{gray}{Leonardo Maccari, \\ leonardo.maccari@unitn.it \\ @leobowski}
	\vspace{.5em}
	\Comment{\\Firenze \\ 19/12/2016}
\end{center}
\end{frame}
}



\begin{frame}[c]{Who am I}
  \bi
  \ii Ricercatore al DISI: Universita' di Trento
  \ii WP coordinator del progetto netCommons: www.netcommons.eu
  \ii Fondatore della rete ninux di Firenze
  \ii @leobowski
  \ei
\end{frame}

%
%\frame{
%    \frametitle{netcommons.eu}
%    \bi
%    \ii Un progetto di ricerca di tre anni, iniziato a Gennaio in cui ci
%    occupiamo di reti comunitarie con un approccio interdisciplinare:
%    \bi
%    \ii 3 dipartimenti di informatica 
%    \ii 2 dipartimenti/centri di ricerca di legge
%    \ii 1 dipartimento di sociologia
%    \ii 1 no-profit
%    \ei
%    \ii Se vi interessa approfondire: www.netcommons.eu
%    \ei
%    %\bab{Posizioni aperte}
%    %\bi
%    %\ii Un dottorato
%    %\ii Un post-doc
%    %\ei
%    %\eab
%}



%%%%%%%%%%%%%%%%%%%%%%%%
% Actual content below %
%%%%%%%%%%%%%%%%%%%%%%%%

\pgfdeclareimage[width=\columnwidth]{prism}{images/PRISM1}
\pgfdeclareimage[width=0.9\columnwidth]{ninux1}{images/ninux1}
\pgfdeclareimage[width=0.6\columnwidth]{ninux2}{images/ninux2}
\pgfdeclareimage[width=0.9\columnwidth]{ninux3}{images/ninux3}
\pgfdeclareimage[width=0.9\columnwidth]{ninux4}{images/ninux4}
\pgfdeclareimage[width=0.9\columnwidth]{ninux5}{images/ninux5}
\pgfdeclareimage[width=0.9\columnwidth]{awmn}{images/awmn}
\pgfdeclareimage[width=0.9\columnwidth]{guifi}{images/guifi}

\frame{
    \frametitle{PRISM}
    \begin{center}
    \pgfuseimage{prism}
    \end{center}
}
\frame{
    \frametitle{Reti Comunitarie}
\bi 
\ii Reti decentralizzate, fatte da persone, per sovvertire questo stato di cose a partire dalla rete stessa
\ei
}


\frame{
    \frametitle{Nodo Ninux}
    \begin{center}
    \pgfuseimage{ninux1}
    \end{center}
}

\frame{
    \frametitle{Nodo Ninux}
    \begin{center}
    \pgfuseimage{ninux2}
    \end{center}
}

\frame{
    \frametitle{Nodo Ninux}
    \begin{center}
    \pgfuseimage{ninux3}
    \end{center}
}

\frame{
    \frametitle{Nodo Ninux}
    \begin{center}
    \pgfuseimage{ninux4}
    \end{center}
}

\frame{
    \frametitle{Ninux}
    \begin{center}
    \pgfuseimage{ninux5}
    \end{center}
}

\frame{
    \frametitle{AWMN}
    \begin{center}
    \pgfuseimage{awmn}
    \end{center}
}


\frame{
    \frametitle{Guifi.net}
    \begin{center}
    \pgfuseimage{guifi}
    \end{center}
}

\frame{
  \frametitle{Ninux Firenze}
  \centering\includegraphics[width=0.4\columnwidth]{images/ninux-fi}
}


\frame{
  \frametitle{Come e' fatta Firenze}
  \centering\includegraphics[width=0.7\columnwidth]{images/firenze-rilievo}
}

\frame{
  \frametitle{Come e' fatta Firenze - II}
  \includegraphics[width=\columnwidth]{images/panorama}
}

\frame{
  \frametitle{Come e' fatta Firenze - III}
  \bi
  \ii Firenze e' circondata da colline mentre la citta' e' abbastanza piatta e gli
  edifici non troppo alti
  \ii Da un lato non ci sono colline che ostacolano la linea di vista
  \ii Da un altro e' difficile trovare punti alti in cui superare gli ostacoli 
  \ei
}

\frame{
  \frametitle{Cosa ci facciamo?}
  \bi
\ii \ldots lo vedremo nei prossimi talk
  \ei
}


\frame{
  \frametitle{Conclusioni}
  \bi
  \ii Le reti comunitarie sono infrastrutture di rete costruite dal basso, 
  e gestite come beni comuni
  \ii Sono un ottimo modo per sperimentare, imparare, e per riprendere il controllo delle nostre comunicazioni
  \ii Partecipate!
  \bi
\ii www.firenze.ninux.org
\ii firenze@ninux.org 
  \ei
  \ei
}



%\frame{
%    \frametitle{netcommons.eu}
%    \bi
%    \ii Un progetto di ricerca di tre anni, iniziato a Gennaio in cui ci
%    occupiamo di reti comunitarie con un approccio interdisciplinare:
%    \bi
%    \ii 3 dipartimenti di informatica 
%    \ii 2 dipartimenti/centri di ricerca di legge
%    \ii 1 dipartimento di sociologia
%    \ii 1 no-profit
%    \ei
%    \ii Se vi interessa approfondire: www.netcommons.eu
%    \ei
%    \bab{Posizioni aperte}
%    \bi
%    \ii Un dottorato
%    \ii Un post-doc
%    \ei
%    \eab
%}
%





\end{document}

